% !TeX root = ../../Book1.tex
% 纲要标题：### 11.4 唯一分解整环
% 纲要位置：计划纲要.md，第 434 行。

\section{唯一分解整环}
\label{b1:sec:1104}

% 纲要内容（仅作写作计划，尚非教材正文）：
%
% * 分解的存在和唯一性
% * 唯一分解不蕴含主性
% * 非唯一分解例子与理想分解的动机
%
% ---
%

主理想性给出了唯一分解的充分条件，却不是必要条件。本节把分解论自身的要求提炼出来，并考察两个相反方向的例子：\(\mathbb Z[x]\) 的元素仍能唯一分解，但有非主理想；\(\mathbb Z[\sqrt{-5}]\) 连元素的分解唯一性也会失败。前者的唯一分解性将在下一章通过多项式方法证明，后者在这里直接计算。

\subsection{分解的存在性与唯一性}
\label{b1:subsec:1104:01}

\begin{proposition}\label{b1:prop:ufd-prime}
设 \(R\) 为整环。如果 \(R\) 是唯一分解整环，则 \(R\) 中的不可约元都是素元。更一般地，\(R\) 是唯一分解整环，当且仅当它的每个非零非单位都能分解成有限个不可约元的乘积，并且每个不可约元都是素元。
\end{proposition}
\begin{proof}
先设 \(R\) 为唯一分解整环，\(p\) 不可约且 \(p\mid ab\)。若 \(a=0\) 或 \(b=0\)，结论成立；否则写 \(ab=pc\)，其中 \(c\ne0\)。把 \(a,b,c\) 分别分解为单位与不可约元的乘积，允许单位的分解含零个不可约因子。等式右侧有因子 \(p\)，唯一性要求它与左侧某个不可约因子相伴。那个因子来自 \(a\) 或 \(b\)，所以 \(p\mid a\) 或 \(p\mid b\)，即 \(p\) 为素元。

反过来，假定分解存在且不可约元都为素元。\cref{b1:thm:pid-ufd}证明唯一性的后一段，在已有两种分解后只用了不可约元的素性、整环的非零消去律及单位的基本性质，并未再用理想主性。因此同一因子配对论证适用于此，得到因子个数相同且对应因子相伴，正是\cref{b1:def:ufd}的唯一性。
\end{proof}

这个判别不能省略分解存在的条件。证明唯一性时以“取一个不可约分解”起步，只在确知每个非零非单位都能如此分解时才有意义。\cref{b1:thm:pid-ufd}为主理想整环补足存在性时，用到的是理想升链最终稳定，两个论证的职责应当分清。

\begin{proposition}\label{b1:prop:ufd-gcd-lcm}
设 \(R\) 为唯一分解整环，\(a,b\in R\) 非零。则 \(a,b\) 有最大公因子。若把它们分解中涉及的不同不可约相伴类各选一个代表 \(p_1,\ldots,p_t\)，其中 \(t\geq0\) 为整数，并写成
\[
a=u\prod_{j=1}^{t}p_j^{\alpha_j},\qquad
b=v\prod_{j=1}^{t}p_j^{\beta_j},\qquad \alpha_j,\beta_j\in\mathbb Z_{\geq0},
\]
其中 \(u,v\in R^\times\)，则可取
\[
d=\prod_{j=1}^{t}p_j^{\min(\alpha_j,\beta_j)},\qquad
m=\prod_{j=1}^{t}p_j^{\max(\alpha_j,\beta_j)}.
\]
其中 \(d\) 为最大公因子；\(m\) 同时被 \(a,b\) 整除，并整除它们的每个公倍数。
\end{proposition}
\begin{proof}
唯一分解说明，对于非零元素 \(r,s\)，条件 \(r\mid s\) 等价于 \(r\) 中每个不可约因子的指数不超过它在 \(s\) 中的指数：必要性由 \(s=rc\) 及分解唯一性得到，充分性由指数相减构造商得到。因此公因子中各指数的上界是两个指数的较小者，公倍数中各指数的下界是两个指数的较大者，分别给出所列 \(d,m\)。公倍数为 \(0\) 时，\(m\mid0\) 也成立。
\end{proof}

设 \(a,b\) 为整环 \(R\) 的非零元素。如果 \(m\in R\) 同时被 \(a,b\) 整除，并且 \(m\) 整除 \(a,b\) 的每个公倍数，就称 \(m\) 为 \(a,b\) 的\term[zuixiaogongbeishu]{最小公倍数}[least common multiple]，记为 \(\operatorname{lcm}(a,b)\)。最小公倍数只在相伴意义下确定。对上一命题选出的 \(d,m\)，乘积 \(dm\) 与 \(ab\) 相伴；适当调整单位后可令二者相等。零元素的情况另定为 \(\gcd(a,0)\) 与 \(a\) 相伴、\(\operatorname{lcm}(a,0)=0\)。
\symidx{\(\operatorname{lcm}(a,b)\)：最小公倍数}

\ProseParagraphBreak
例如在 \(\mathbb Z[i]\) 中，\(2=-i(1+i)^2\)，所以 \(2\) 已不再是不可约元。若 \(a=(1+i)^3(2+i)\)、\(b=(1+i)(2+i)^2\)，则 \(1+i\) 与 \(2+i\) 的范数分别为整数素数 \(2,5\)，两者不可约且不相伴。因而可取 \(\gcd(a,b)=(1+i)(2+i)\)，最小公倍数则为 \((1+i)^3(2+i)^2\)。

\subsection{唯一分解不蕴含主性}
\label{b1:subsec:1104:02}

\cref{b1:prop:integer-poly-nonprincipal}已经说明：\(\mathbb Z[x]\) 的理想 \((2,x)\) 不为主理想。这里的障碍是 \(2,x\) 虽然没有共同的不可约因子，却不能用多项式系数的线性组合得到 \(1\)。

\ProseParagraphBreak
另一方面，\(\mathbb Z[x]\) 确实是唯一分解整环。这个断言是下一章\cref{b1:thm:polynomial-ufd}的一个应用：那里将证明只要 \(R\) 是唯一分解整环，\(R[x]\) 也是；取本章已证的 \(R=\mathbb Z\) 即得。此处把它作为前瞻例子，不用它证明本章其他结论。与已经完成证明的 \(\mathbb Z[\omega]\) 一起，这两个例子说明三个条件之间是严格蕴含：
\[
\text{Euclid 整环}\ \Longrightarrow\
\text{主理想整环}\ \Longrightarrow\
\text{唯一分解整环}.
\]
左边逆向的反例由\cref{b1:prop:noneuclidean-pid}给出，右边逆向的反例则在下一章完成最后的唯一分解性验证。

\ProseParagraphBreak
还有一个可以立即比较的区别。在唯一分解整环中，对非零 \(a,b\)，公倍数恰为最小公倍数的倍数，因此 \((a)\cap(b)=(m)\) 仍是主理想；而和 \((a)+(b)\) 未必为主理想。最大公因子刻画公因子之间的整除关系，却未必能表示为原来两个元素的线性组合。将“有最大公因子”与“能写 Bézout 关系”分开，能避免把主理想整环的结论错误推广到全部唯一分解整环。

\subsection{非唯一分解与理想分解}
\label{b1:subsec:1104:03}

令 \(s=\sqrt{-5}\)，考虑整环 \(R=\mathbb Z[s]\subseteq\mathbb C\)。它的元素唯一写成 \(a+bs\)，范数
\[
N(a+bs)=a^2+5b^2
\]
是乘法相容的非负整数，单位恰为 \(\pm1\)。不存在范数为 \(2\) 或 \(3\) 的元素：若 \(b\ne0\)，范数至少为 \(5\)；若 \(b=0\)，整数平方又不可能等于 \(2,3\)。

\begin{proposition}\label{b1:prop:sqrt-minus-five-factorization}
令 \(s=\sqrt{-5}\)，并令 \(R=\mathbb Z[s]\subseteq\mathbb C\)。则等式
\[
6=2\cdot3=(1+s)(1-s)
\]
给出两种不等价的不可约分解。因此 \(R\) 不是唯一分解整环，且不可约元 \(2\) 不是素元。
\end{proposition}
\begin{proof}
元素 \(2,3,1+s,1-s\) 的范数依次为 \(4,9,6,6\)。若 \(2\) 分解成两个非单位，两因子的范数只能为 \(2,2\)，但范数 \(2\) 不存在，故 \(2\) 不可约。对 \(3\)，唯一的非平凡范数分解为 \(3\cdot3\)，同样不可能。对 \(1\pm s\)，非平凡范数分解只能为 \(2\cdot3\)，也不可能。因此四个因子都不可约。

单位的范数为 \(1\)，相伴元素范数相同，所以范数为 \(4\) 的 \(2\) 不可能与范数为 \(6\) 的任一 \(1\pm s\) 相伴。两种分解因而不等价。此外，\(2\mid(1+s)(1-s)=6\)，却不整除任一 \(1\pm s\)，因为它们的整数常数项为奇数。故 \(2\) 不满足素性。
\end{proof}

元素分解失效以后，可以尝试让理想承担因子的作用。本例的
\[
I=(2,1+s)
\]
不是主理想，却满足 \(I^2=(2)\)。先看非主性。映射
\(R\rightarrow\mathbb F_2\)、\(a+bs\mapsto\overline{a+b}\)
保持乘法，因为 \(s^2=-5\) 在模 \(2\) 下对应 \(1^2=1\)；它把 \(I\) 送到零而把 \(1\) 送到 \(1\)，故 \(I\ne R\)。若 \(I=(d)\)，则 \(d\mid2\)。由 \(2\) 不可约，\(d\) 或为单位，或与 \(2\) 相伴；前者与 \(I\ne R\) 矛盾，后者又使 \(1+s\in(2)\)，也矛盾。

\ProseParagraphBreak
再算理想平方。\(I^2\) 由
\[
4,\quad 2+2s,\quad (1+s)^2=-4+2s
\]
生成，三者都属于 \((2)\)。反向包含由
\[
2=(2+2s)-(-4+2s)-4\in I^2
\]
得到，所以 \(I^2=(2)\)。元素 \(2\) 不可约，但它生成的理想仍有这样的真分解。\cref{b1:ex:11-ideal-splitting-three}将继续核对 \(J=(3,1+s)\)、\(K=(3,1-s)\) 满足 \(JK=(3)\)、\(IJ=(1+s)\)、\(IK=(1-s)\)。于是原来的两种元素分解在理想层面汇合为
\[
(6)=(2)(3)=I^2JK=(IJ)(IK)=(1+s)(1-s).
\]
最后一项表示两个主理想的乘积。这些等式只核对当前例子；要得到一般的理想分解定理，还须对环施加适当条件。这是后续数论与交换代数需要研究的问题。
